%! Equations and Graphs in Both Directions — Unit 2, Lesson 2.3 — Warm-Up (Answer Key)
\documentclass[10pt]{article}
\usepackage{atda-article}
\usepackage{atda-key}   % pulls in -boxes; do NOT also load -boxes
\newcommand{\TallMath}[1]{$\displaystyle #1\rule[-1.4em]{0pt}{3.2em}$}
\setlength{\workrowsep}{6pt}   % handwriting room; moves blank and key together

\begin{document}

\pageheader{Unit 2, Lesson 2.3}{Warm-Up --- Answer Key}
% NAMESTRIP: no \namedateperiod here — mirrors the blank. Teacher-only prose goes
% in the lesson plan as \begin{teachernote}[Warm-Up], never in a key.

\vspace{0.15in}

\begin{notesbox}{Same Keys, Different Meaning}
\small
Five quick items. Item 2 is two expressions that are typed almost the same way and are not the same
rule --- read it slowly.

\begin{enumerate}[leftmargin=1.6em, itemsep=4pt, topsep=4pt]

\item \textbf{One rule, one table.} Let $f(x)=(x-2)^2$. Fill in every cell.
\smallskip

\renewcommand{\arraystretch}{1.3}
\begin{tabular}{@{}|c|c|c|c|c|c|@{}}
\hline
$x$ & $0$ & $1$ & $2$ & $3$ & $4$ \\ \hline
$f(x)=(x-2)^2$ & \ans{$4$} & \ans{$1$} & \ans{$0$} & \ans{$1$} & \ans{$4$} \\
\hline
\end{tabular}
\smallskip

Which input gives the smallest output, and what is that output? \ans{$x=2$ gives $f(2)=0$.}

\item \textbf{Where the parentheses go.} Evaluate both of these at $x=5$.
\begin{work}
  2^x-5 &= 2^5-5 \\
        &= 32-5 \\
        &= 27 \\
  2^{x-5} &= 2^{5-5} \\
          &= 2^0 \\
          &= 1
\end{work}

\item \textbf{Naming the move} (Lesson 2.2). Name the parent function, then describe the
transformation. \quad (a) $y=(x+4)^2$ \quad (b) $y=2^x-6$ \quad (c) $y=-x^2$

\ansline{(a) Parent $f(x)=x^2$, translated \textbf{left $4$} (the $4$ is inside the parentheses).}\\
\ansline{(b) Parent $h(x)=2^x$, translated \textbf{down $6$}. \quad (c) Parent $f(x)=x^2$, reflected over the $x$-axis.}\\

\item \textbf{Writing a set down} (Lesson 2.1). Give the domain and the range of $y=x^2+3$ in
interval notation.

\ansline{Domain $(-\infty,\infty)$; range $[3,\infty)$ --- a bracket, because $3$ is reached at $x=0$.}\\

\item \textbf{Last night's homework, in one sentence.} You wrote an equation for the parabola $y_2$
and were asked what a graphing calculator would show if your equation was right. What did you
predict?

\ansline{$y_2=(x+2)^2$: the screen should show the parent parabola shifted $2$ units left, lowest point $(-2,0)$.}\\

\end{enumerate}
\end{notesbox}

\end{document}
